\documentclass[14pt,aspectratio=1610]{beamer}
%\documentclass[14pt,aspectratio=1610,handout]{beamer}
\usepackage{csu-theme}
% \usepackage{linked-lists}
\usepackage{binary-trees}
\usepackage{booktabs}

\forestset{
  default preamble={%
    for tree={%
      treeNode,%
      %font=\small,%
      % l sep=0.2em,%
      % level distance=0.5em,%
      % outer sep=0pt,
      % outer ysep=0pt,
      l=2.4em,
    },%
    for leaves={%
      justText},
  },%
}

\title{\Huge{} Binary Search Trees\\Deletion \& Generic Traversal}
\subtitle{CSCI 315: Data Structure Analysis\\Chapter 19}
%\author{Sean T. Hayes, Ph.D.}
\institute{Charleston Southern University}

%% Comment out date to use default (today)
% \date{March 6, 2025}
\subject{Software Programming}
%\keywords{}
\graphicspath{{../imgs/binary-tree}}

\setlength{\parskip}{0.5em}

\begin{document}

\begin{frame}
  \titlepage{}
\end{frame}

\section{Review}
\begin{frame}{Binary Trees}
\begin{columns}
  \column{0.3\textwidth}
    A \emph{binary tree} is empty or has these properties:
    \begin{itemize}
    \item
      Has a root node.
    \item
      Has a left subtree and right subtree, which are binary trees.
    \end{itemize}
  \column{0.62\textwidth}
    \centering
    \begin{forest}
    [root,headPtr, for tree={fit=band, s sep=6pt}
      [A
        [B
          [D,
            [\NullNode]
            [G
              [\NullNode]
              [\NullNode]
            ]
          ]
          [E
            [\NullNode]
            [\NullNode]
          ]
        ]
        [C
          [\NullNode]
          [F
            [H
              [\NullNode]
              [\NullNode]
            ]
            [\NullNode]
          ]
        ]
      ]
    ]
    \end{forest}
\end{columns}
\end{frame}

\begin{frame}{Binary Trees}
  \begin{columns}
  \column{0.55\textwidth}
    \begin{itemize}
      \item
        A node (called a vertex in graph theory):
        \begin{itemize}
        \item
          Has at most two children.
        \item
          Stores its own information.
        \item
          Keeps track of its left subtree and right subtree using pointers\\
          (e.g., \lstinline{left} and \lstinline{right}).
        \end{itemize}
      \item
         A pointer to the root node of the binary tree is stored outside the tree.
    \end{itemize}
  \column{0.45\textwidth}
    \def\svgwidth{\textwidth}
    \input{../imgs/binary-tree/binary-tree-linked.pdf_tex}
    %\includegraphics[width=\textwidth]{binary-tree-linked}
  \end{columns}
\end{frame}

\begin{frame}{Binary Trees}
  \begin{columns}
    \column{0.5\textwidth}
      \begin{itemize}[<+->]
      \item
        \emph{Leaf}: a node that has no children.
      \item
        A \emph{parent} has a branch to a \emph{child} node.
      \end{itemize}
    \column{0.5\textwidth}
      \centering
      \begin{forest}
      [A, for tree={fit=tight}
        [B
          [D,
            [\NullNode]
            [G
              [\NullNode]
              [\NullNode]
            ]
          ]
          [E
            [\NullNode]
            [\NullNode]
          ]
        ]
        [C
          [\NullNode]
          [F
            [H
              [\NullNode]
              [\NullNode]
            ]
            [\NullNode]
          ]
        ]
      ]
      \end{forest}
  \end{columns}
\end{frame}

\begin{frame}{Binary Trees}
  \begin{columns}
    \column{0.5\textwidth}
      \begin{itemize}[<+->]
      \item
        There is a unique \emph{path} from the root to every node.
      \item
        The \emph{Path Length} is the number of branches on that path.
      \end{itemize}
    \column{0.5\textwidth}
      \centering
      \begin{forest}
      [A, for tree={fit=tight}
        [B
          [D,
            [\NullNode]
            [G
              [\NullNode]
              [\NullNode]
            ]
          ]
          [E
            [\NullNode]
            [\NullNode]
          ]
        ]
        [C
          [\NullNode]
          [F
            [H
              [\NullNode]
              [\NullNode]
            ]
            [\NullNode]
          ]
        ]
      ]
      \end{forest}
  \end{columns}
\end{frame}
  
\begin{frame}{Binary Trees}
  \begin{columns}
    \column{0.5\textwidth}
      \begin{itemize}[<+->]
      \item
        Node \emph{Level}: number of branches from the root to the node.
        \begin{itemize}
        \item
          The root node is at level 0.
        \end{itemize}
      \item
        \emph{Height} of a binary tree: number of edges on the longest path from the root to a leaf.
      \end{itemize}
    \column{0.5\textwidth}
      \centering
      \begin{forest}
      [A, for tree={fit=tight}
        [B
          [D,
            [\NullNode]
            [G
              [\NullNode]
              [\NullNode]
            ]
          ]
          [E
            [\NullNode]
            [\NullNode]
          ]
        ]
        [C
          [\NullNode]
          [F
            [H
              [\NullNode]
              [\NullNode]
            ]
            [\NullNode]
          ]
        ]
      ]
      \end{forest}
  \end{columns}
\end{frame}

\begin{frame}{Binary Search Tree}

  \begin{columns}
  \column{0.60\textwidth}
    Definition: A \emph{binary search tree} is either empty or has these properties:
    \begin{itemize}
      \item
        Has a root node.
      \item
        Has two sets of nodes:\\the left and right subtrees.
      \item
        The key in the root node is larger than every key in the left subtree and smaller than every key in the right subtree.
      \item
        The left subtree and right subtree are binary search trees.
    \end{itemize}
  \column{0.29\textwidth}
    \begin{forest}
      [96, for tree={fit=tight}
        [27
          [20
            [\NullNode]
            [\NullNode]
          ]
          [45
            [30
              [\NullNode]
              [\NullNode]
            ]
            [\NullNode]
          ]
        ]
        [\NullNode]
      ]
    \end{forest}
  \end{columns}
\end{frame}
  
\begin{frame}{Performance of a Binary Search Tree}
  What is the \textit{average-case} Big-O of the following operations on a binary search tree?\\[1em]
  
  \renewcommand{\thefootnote}{\fnsymbol{footnote}}
  \begin{tabular}{r c c c}
  \toprule
                  & \textbf{Array} & \textbf{Linked List} & \textbf{Binary Search Tree}\\
  \midrule
    Random Access & {\(O(1)\)} & \(O(n)\) & {\textit{NA}}\\
    Insert/Remove & {\(O(n)\)} & \(O(n)\) & {\(O\left(\log_2{n}\right)\)}\\
    Search        & {\(O(n)\)} & \(O(n)\) & {\(O\left(\log_2{n}\right)\)}\\
    Sort          & {\(O\left(n \log_2{n}\right)\)} & {\(O\left(n \log_2{n}\right)\)} & {\textit{NA}\footnote[frame]{{\small A binary-search tree is already sorted. To convert it to a \\sorted list, a inorder traversal will take \(O(n)\).\\[1em]}}}\\
  \bottomrule
  \end{tabular}
\end{frame}

\section{Delete}
\begin{frame}{Deleting a Node}
\begin{columns}
% \setbeamercovered{transparent}
\column{0.6\textwidth}
  %\begin{center}
    \begin{forest}
      [root, headPtr, for tree={fit=tight, l=2.4em}, for leaves={justText},%l sep?
        [60
          [\alt<-9>{50}{40},name=delType4
            [30,
              [\NullNode]
              [35
                [32
                  [\NullNode]
                  [\NullNode]
                ]
                [\alt<-9>{40}{\NullNode}, name=delType4Leaf, justText on=<10->, for descendants={visible on=<-9>}
                  [\NullNode, visible on=<-9>]
                  [\NullNode, visible on=<-9>]
                ]
              ]
            ]
            [53, node visible on=<-4> %\alt<-4>{53}{57}
              [\NullNode, node visible on=<-4>]
              [57, name=delType2Child%\alt<-4>{57}{\NullNode}, justText on=<5->, for descendants={visible on=<-4>},
                [\NullNode]
                [\NullNode]
              ]
            ]
          ]
          [70, name=delType4Parent
            [\NullNode]
            [80, node visible on=<-6> %\alt<-6>{80}{75}
              [75, name=delType4Child %\alt<-6>{75}{\NullNode}, justText on=<7->, for descendants={visible on=<-6>}
                [\NullNode]
                [\alt<-2>{77}{\NullNode}, justText on=<3->, for descendants={visible on=<-2>}
                  [\NullNode]
                  [\NullNode]
                ]
              ]
              [\NullNode, node visible on=<-6>]
            ]
          ]
        ]
      ]
      \onslide<5->{
        \draw[->, very thick] (delType4) to[bend left=20] node[midway,left]{} (delType2Child);
      }
      \onslide<7->{
        \draw[->, very thick] (delType4Parent) to[bend left] node[midway,left]{} (delType4Child);
      }
      \onslide<9>{
        \draw[->,dotted, very thick] (delType4Leaf) to[bend right=15] node[midway,left] {\small{}move} (delType4);
      }
      \only<10>{}
    \end{forest}
%\end{center}

\column{0.4\textwidth}
  The four cases of delete depend on the node to be deleted\textellipsis
  \begin{enumerate}
    \item<2-> Is a leaf (e.g., 77)
    \item<4-> Has an empty left subtree (e.g., 53)
    \item<6-> Has an empty right subtree (e.g., 80)
    \item<8-> Has nonempty left and right subtrees (e.g., 50)
  \end{enumerate}
\end{columns}
\end{frame}


\begin{frame}{Example Tree}
%\begin{columns}
% \setbeamercovered{transparent}
%\column{0.6\textwidth}
  %\begin{center}
  \centering{}
    % \forestset{
    %   treeNode/.append style={inner sep=2pt}
    % }
    \begin{forest}
      for tree={s sep=1mm}
      [root,headPtr%, l sep=0mm, for tree={fit=tight, l sep=-5mm}, for leaves={l sep=0mm}
        [60,l sep=0mm
          [50
            [30
              [\NullNode]
              [35
                [32
                  [\NullNode]
                  [\NullNode]
                ]
                [40
                  [\NullNode]
                  [\NullNode]
                ]
              ]
            ]
            [53
              [\NullNode]
              [57
                [\NullNode]
                [\NullNode]
              ]
            ]
          ]
          [70
            [\NullNode]
            [80
              [75
                [\NullNode]
                [77
                  [\NullNode]
                  [\NullNode]
                ]
              ]
              [\NullNode]
            ]
          ]
        ]
      ]
    \end{forest}
%\end{center}
\end{frame}

%% To Do: Add an example of removing a node with two children.

\begin{frame}{Example Tree: Remove 50}
%\begin{columns}
%\column{0.6\textwidth}
  %\begin{center}
    % \forestset{
    %   treeNode/.append style={inner sep=2pt}
    % }
    %\centering{}
    \begin{forest}
      for tree={s sep=1mm}
      %[root,headPtr, l sep=0mm
        [60,%l sep=0mm, for tree={fit=tight, l sep=-5mm}, for leaves={l sep=0mm}
          [50, name=el50, edge highlight on=<2->
            [30, %edge highlight on=<2>
              [\NullNode]
              [35, name=el35%, edge highlight on=<2>
                [32
                  [31
                    [\NullNode]
                    [\NullNode]
                  ]
                  [\NullNode]
                ]
                [40, name=el40, edge highlight on=<3->, edge+={/tikz/visible on=<-3>}
                  [39, name=el39
                    [\NullNode]
                    [\NullNode]
                  ]
                  [\NullNode]
                ]
              ]
            ]
            [53
              [\NullNode]
              [57
                [\NullNode]
                [\NullNode]
              ]
            ]
          ]
          [70
            [\NullNode]
            [80
              [75
                [\NullNode]
                [77
                  [\NullNode]
                  [\NullNode]
                ]
              ]
              [\NullNode]
            ]
          ]
        ]
      \begin{scope}[->,very thick,white]
        % \draw[->,dotted, very thick] () to[bend left] node[midway,left] {rotate} (p2);
        \onslide<2->{
          \node[anchor=east, inner sep=1pt] (txtLabel) at ([shift={(-0.4,0.4)}]el50.north west) {\texttt{\small pREMOVE}};
          \draw(txtLabel.east) to[bend left] (el50);
        }
        \onslide<3->{
          \node[anchor=west, inner sep=1pt] (txtLabel) at ([shift={(0.2,0.45)}]el35.north east) {\texttt{\small pToMoveParent}};
          \draw(txtLabel.west) to[bend right] (el35);
          \node[anchor=west, inner sep=1pt] (txtLabel) at ([shift={(0.2,0.25)}]el40.north east) {\texttt{\small pToMove}};
          \draw(txtLabel.west) to[bend right] (el40);
        }
        \onslide<4->{
          \draw[highlight] (el35) to[bend left] (el39);
        }
      \end{scope}
    \end{forest}
%\end{center}
\end{frame}


\begin{frame}{Example Tree: Remove 35}
  %\begin{columns}
  %\column{0.6\textwidth}
    %\begin{center}
      % \forestset{
      %   treeNode/.append style={inner sep=2pt}
      % }
      %\centering{}
      \begin{forest}
        for tree={s sep=1mm}
        %[root,headPtr, l sep=0mm
          [60,%l sep=0mm, for tree={fit=tight, l sep=-5mm}, for leaves={l sep=0mm}
            [50, name=el50
              [30, %edge highlight on=<2>
                [\NullNode]
                [35, name=el35, edge highlight on=<2->
                  [32, name=el32, edge highlight on=<3->, edge+={/tikz/visible on=<-3>}
                    [31, name=el31
                      [\NullNode]
                      [\NullNode]
                    ]
                    [\NullNode]
                  ]
                  [40, name=el40
                    [39, name=el39
                      [\NullNode]
                      [\NullNode]
                    ]
                    [\NullNode]
                  ]
                ]
              ]
              [53
                [\NullNode]
                [57
                  [\NullNode]
                  [\NullNode]
                ]
              ]
            ]
            [70
              [\NullNode]
              [80
                [75
                  [\NullNode]
                  [77
                    [\NullNode]
                    [\NullNode]
                  ]
                ]
                [\NullNode]
              ]
            ]
          ]
        %]
              % \node[] at (el50.north west) (deletePtr) {toDelete};
        \begin{scope}[->,very thick,white]
          % \draw[->,dotted, very thick] () to[bend left] node[midway,left] {rotate} (p2);
          \onslide<2->{
            \node[anchor=west, inner sep=1pt] (txtLabel) at ([shift={(0.2,0.45)}]el35.north east) {\texttt{\small pREMOVE}};
            \draw(txtLabel.west) to[bend right] (el35);
          }
          \onslide<3->{
            % \node[anchor=west, inner sep=1pt] (txtLabel) at ([shift={(0.2,0.45)}]el35.north east) {\texttt{\small pToMoveParent}};
            % \draw(txtLabel.west) to[bend right] (el35);
            \node[anchor=east, inner sep=1pt] (txtLabel) at ([shift={(-0.2,0.45)}]el32.north west) {\texttt{\small pToMove}};
            \draw(txtLabel.east) to[bend left] (el32);
          }
          \onslide<4->{
            \draw[highlight] (el35) to[bend left] (el31);
          }
        \end{scope}
      \end{forest}
  %\end{center}
\end{frame}

\begin{frame}{Performance of Delete}
  \centering{}\Large
  What is the average-case Big-\(O\) performance of delete?
\end{frame}

\section{Generic Traversal}
\begin{frame}{Generic Traversal}
  \begin{itemize}[<+->]
    \item
      In a traversal algorithm, ``visiting'' may mean different things.
      \begin{itemize}
        \item For example, outputting the values, updating the values, and calculating results.
      \end{itemize}
    \item
      \emph{Problems}:
      \begin{itemize}
        \item
          Repeatedly implementing a specific traversal algorithm for each function.
        \item 
          Must have access to the internals of the tree.
      \end{itemize}
    \item
      \emph{The Solution}: Write a generic traversal function that accepts a function as a parameter.
  \end{itemize}
\end{frame}

\begin{frame}{Generic Traversal}
  \begin{itemize}[<+->]
    \item
      In C++, a function name without parentheses is considered a pointer to the function.
    \item
      To specify a function as a formal parameter to another function,
      \begin{enumerate}
        \item specify the function type (e.g., \lstinline{void}),
        \item the name as a pointer (e.g., \lstinline{(*visit)}),
        \item the parameter types (e.g., \lstinline{(int, double)}).
      \end{enumerate}
    \item For example,\\
      \lstinline{void funcParamFunc1(void (*visit) (int));}\\[1em]
      \lstinline{template<typename Type>}\\
      \lstinline{void funcParamFunc2(void (*visit) (Type&));}
  \end{itemize}
\end{frame}

\begin{frame}[fragile]{Generic Traversal Example}
A recursive member function that accepts a node to tree to traverse and a function to call at each node.
\begin{lstlisting}[basicstyle=\ttfamily\small]
template<typename Type>
void BinarySearchTree<Type>::inorder(Node *pRoot,
      void (*visit)(Type& value)) {
  if (pRoot != nullptr)  {
    inorder(pRoot->left, visit);

    visit(pRoot->value);

    inorder(pRoot->right, visit);
  }
}
\end{lstlisting}
\end{frame}

\section{Function Objects}
\begin{frame}{Function Objects (a.k.a. Functors)}
  \begin{itemize}[<+->]
    \item
      In short, and object that looks like a function, but that\\
      can maintain state.
    \item
      A function object is an object class/struct that overloads \lstinline{operator()}.
    \item
      Can be passed to a generic traversal function instead of a function pointer.
  \end{itemize}
\end{frame}

\begin{frame}[fragile]{Function Object: Example Definition}
\begin{lstlisting}[basicstyle=\ttfamily\small]
template <typename Type>
struct SumValues {
  Type sum;

  // Constructor initializes sum to 0
  SumValues() : sum(0) {}

  // Overloaded operator() to accumulate sum
  void operator()(Type data) {
    sum += data;
  }
};
\end{lstlisting}
\end{frame}

\begin{frame}[fragile]{Function Object: How it Works}
\begin{lstlisting}[basicstyle=\ttfamily\small]
SumValues funObj;

funObj(10);
funObj(20);

std::cout << funObj.sum << '\n'; // 30
\end{lstlisting}
\end{frame}

\begin{frame}[fragile]{Function Object: Example Generic Traversal}
\begin{lstlisting}[basicstyle=\ttfamily\small]
template<typename Type>
template<typename Func>
void BinarySearchTree<Type>::preorderTraversal( Node* pCurr, Func& functor)
{
  if (pCurr != nullptr)
  {
    functor(pCurr->value); // visit the node
    preorderTraversal(pCurr->pLeft, functor);
    preorderTraversal(pCurr->pRight, functor);
  }
}
\end{lstlisting}
\end{frame}

\begin{frame}[fragile]{Function Object: Example Generic Traversal}
\begin{lstlisting}[basicstyle=\ttfamily\small]
template<typename Type>
template<typename Func>
void BinarySearchTree<Type>::preorderTraversal( Func& functor)
{
  preorderTraversal(mpRoot, functor);
}
\end{lstlisting}
\end{frame}

\begin{frame}[fragile]{Function Object: Example Generic Traversal}
\begin{lstlisting}[basicstyle=\ttfamily\small]
BinarySearchTree<int> tree{VALUES, SIZE};
SumValues<int> sumObj; // Instantiate functor

tree.preorderTraversal(sumObj);

std::cout << "Sum: " << sumObj.sum << '\n';
\end{lstlisting}
\end{frame}

\section{Self-Balancing}
\begin{frame}{Self-Balancing Binary Search Trees}
\begin{itemize}[<+->]
  \item
    The worst-case performance for searching is \(O(n)\).
  \item
    \emph{Self-Balancing Binary Trees} automatically \\
    minimize the tree's height when performing\\
    insertion and deletion operations.
  \item
    Searching (and therefore insertion and deletion) becomes $O(log_2{n})$ (i.e., asymptotically optimal) in the worst case (and may average $O(1)$ for some operation depending on the implementation).
  \item
    \emph{Disadvantage}: Algorithms are more complicated (especially if duplicates are allowed).
\end{itemize}
\end{frame}

\begin{frame}{Self-Balancing Binary Search Trees}
  Examples:
\begin{itemize}
  \item
    2--3
  \item 
    AA
  \item
    AVL
  \item
    Red Black
  \item
    Scapegoat
  \item
    Splay
  \item
    Treap
  \item
    etc.
\end{itemize}
\end{frame}

\begin{frame}{Example of Balancing using Rotation}
\centering

\begin{forest}
  [,phantom, for tree={fit=band, l sep-=1em, minimum size=1.5em, font=\normalsize}, for leaves={treeNode}, for tree={s sep=0.5cm}, s sep=1.5cm
    [3
      [1, name=p1
        [,justText, no edge, for current={no edge}, name=hid1] {
          \onslide<2->{\draw[<-,dotted, very thick] (hid1) to[bend left] node[midway,left] {rotate} (p1);}
        }
        [2,name=c1] {
          \onslide<2->{\draw[->,dotted, very thick] () to[bend right] node[midway,right] {rotate} (p1);}
        }
      ]
      [, phantom]
    ]
    [3, name=p2, for tree={visible on=<3->}
      [2
        [1]
        [, phantom]
      ]{
        \onslide<4->{\draw[->,dotted, very thick] () to[bend left] node[midway,left] {rotate} (p2);}
      }
      [,justText, for current={no edge}, name=hid2] {
        \onslide<4->{\draw[<-,dotted, very thick] () to[bend right] node[midway,right] {rotate} (p2);}
      }
    ]
    [2, for tree={visible on=<5->}
      [1]
      [3]
    ]
  ]
\end{forest}
\end{frame}

\section*{Lab 15}
\begin{frame}{Lab 15: Binary Trees (Part 2)}
\begin{center}\Large
Let's take a look at the lab based on today's lecture.
\end{center}
\end{frame}
\end{document}
